% Pursuing Stacks (cotes 134-2 to 134-8): archive cote 134-3, site page 33 (PDF page 34), typescript page 33.
% typescript. Draft compared: ps2.tex:1567-1638.
% AI-transcribed by Claude Opus 5.5 (claude-opus-5-5, adaptive thinking, effort high), Message Batches API, batch ps1 (msgbatch_01Rv5hCaNTtUt4JHx1yqoHPL), 2026-11-12; not yet checked by eye.
\scanpage{134-3/33}
\mspage{\hand{33}}
which presumably (as I already felt yesterday, cf first lines p.24) is not
any less model-preserving than the functors $\hand{\alpha},\hand{\beta}$ it generalizes.  Of course,
taking A = $T_{p,q}$ above,  we equally get a natural functor
\[ \hand{i_{p,q}} : T^{\wedge}_{p,q} \longrightarrow (\mathrm{Cat}) \]
\margin{$\uparrow$} admitting a right adjoint\inserted{$j_{p,q}$}, and both functors I fe\unsure{e}l must be model preser-%
ving.

\over{\unsure{26}}{29}.  It is time now to elaborate a little upon the notion of a test catego-%
ry,  within the context of modelizers.  Let A be a small category,  and
consider  the functor (**)
\margin{$\downarrow$} (**) \[ \hand{i}_{A} : A^{\wedge} \longrightarrow (\mathrm{Cat}) \quad , \quad F \mapsto A_{/F} \quad , \]
Whenever we have a functor \inserted{i:} $M \to M'$,  when $M'$ is equipped with a $W'$ turning
it into a modelizer,  there is (if any) \xout{xx} just one $W \subset \mathrm{F}\!\not{\mathrm{l}}(M)$ turning M into
a modelizer and i into a morphism of such,  namely $W = i_{\mathrm{F}\!\not{\mathrm{l}}}^{-1}(W')$.  In any
case,  we may define W (``weak equivalences'') by this formula,  and get a
functor
\[ W^{-1}M \longrightarrow W'^{-1}M' \quad , \]
which is an equivalence iff (M,W) is indeed a modelizer and i model\xout{\unsure{s}}preser-%
ving.  We may say shortly that i: $M \to M'$ is model preserving,  even without
any W given beforehand.  Now coming back to the situation  (**),  the under-%
standing yesterday was to call A a \uline{test} \uline{category} ,  to express that the
canonical functor (**) is model preserving. \hand{(}In any case,  unless otherwise
specified by the context, we will refer to arrows in $A\hand{^{\wedge}}$ which are transfor-%
med into weak equivalences of (Cat)  as ``\uline{weak equivalences}''.\hand{)} It may well
turn out,  by the way,  that we will have to restrict somewhat still the
notion of a test category.

\margin{$\|$} In any case,  \uline{the basic modelizer,}  \uline{in this whole approach to homotopy}
\uline{models,}  \uline{is by no means the category} (Ss sets)  (however handy)  \uline{or the}
\uline{category} (Spaces)  (howeve\xout{\unsure{t}}r appealing to topological intuition),  \uline{but the}
\uline{category} (Cat) \uline{of} \uline{``all''} (small) \uline{categories.}  In this set-up,  the category
(Hot) is most suitably defined as the category of fractions of (Cat) with
respect to ``weak equivalences''.  These in turn are most suitably defined
in cohomological terms,  via the corresponding notion for topoi  -  namely
a morphism of topoi
\[ f:\ X \longrightarrow X' \]
is a ``weak equivalence'' or homotopy equivalence,  iff for every locally
constant sheaf $F'$ on $X'$, the maps
\[ H^{i}(X',F') \to H^{i}(X,f^{-1}(F')) \]
are isomorphisms whenever defined  -  namely for i=0,  for i=1 if moreover
$F'$ is endowed with a group structure,  and for any i if $F'$ is moreover
commutative\xout{t} (\uline{Criterion} \uline{of} \uline{Artin-Mazur}).  Accordingly,  a functor
